% 自写模板 — 华为杯 (中国研究生数学建模竞赛) 论文骨架
% 用法: xelatex main.tex (连编两遍); 或
%       python scripts/render_paper.py --competition huaweibei --workspace <paper_workspace>
% 依据: 研究生数学建模竞赛公开发布的论文格式要求 (题目/摘要/关键词/正文/
%       参考文献/附录的通行结构, 参考文献按 GB/T 7714)。
%       本文件为独立实现, 未参考任何第三方模板源码, 亦未使用竞赛官方
%       视觉资产 (logo 等); 如赛事明确要求官方封面, 请按当年通知自行替换。
% 字体: 正文宋体小四, 章节标题黑体; xelatex 下 ctex 自动选用系统字体。

\documentclass[zihao=-4,a4paper]{ctexart}

% ---------- 页面 ----------
\usepackage{geometry}
\geometry{a4paper, top=2.5cm, bottom=2.5cm, left=2.6cm, right=2.6cm}

% ---------- 数学 ----------
\usepackage{amsmath,amssymb}

% ---------- 图表 ----------
\usepackage{graphicx}
\usepackage{float}
\usepackage{booktabs}
\usepackage{multirow}
\usepackage{longtable}
\usepackage{caption}
\usepackage{subcaption}

% ---------- 代码附录 ----------
\usepackage{listings}
\usepackage{xcolor}

% ---------- 其他 ----------
\usepackage{enumitem}
\usepackage{fancyhdr}
\usepackage[hidelinks]{hyperref}

% 章节标题: 黑体
\ctexset{
  section = {format = {\Large\bfseries\heiti}},
  subsection = {format = {\large\bfseries\heiti}},
  subsubsection = {format = {\bfseries\heiti}},
}

% 图表标题
\captionsetup{font={small,bf}, labelsep=quad}

% 页眉: 左侧赛事名, 右侧页码 (不需要可整段删除, 改用 \pagestyle{plain})
\pagestyle{fancy}
\fancyhf{}
\fancyhead[L]{\small 研究生数学建模论文}
\fancyhead[R]{\small 第 \thepage\ 页}
\renewcommand{\headrulewidth}{0.4pt}

% 代码样式
\lstset{
  basicstyle=\small\ttfamily,
  breaklines=true,
  frame=single,
  numbers=left,
  numberstyle=\tiny\color{gray},
  keywordstyle=\color{blue},
  commentstyle=\color{green!50!black},
  stringstyle=\color{red!60!black},
  showstringspaces=false,
  columns=flexible,
}

\begin{document}

% ============================================================
% 封面 (信息占位区; 按当年参赛要求增删字段)
% ============================================================
\begin{titlepage}
  \begin{center}
    \vspace*{2.0cm}
    {\Huge\heiti [论文标题]}\\[2.5cm]

    {\large
      \begin{tabular}{rl}
        题号:         & [题号] \\[0.8em]
        学校名称:     & [学校名称] \\[0.8em]
        参赛队号:     & [参赛队号] \\[0.8em]
        队员一:       & [队员一姓名] \\[0.8em]
        队员二:       & [队员二姓名] \\[0.8em]
        队员三:       & [队员三姓名] \\[0.8em]
        指导教师:     & [指导教师姓名] \\
      \end{tabular}}
    \vfill
  \end{center}
\end{titlepage}

% 封面不编页码; 页码从摘要页起算
\pagenumbering{arabic}

% ============================================================
% 摘要页
% ============================================================
\begin{center}
  {\Large\heiti 摘\quad 要}
\end{center}

% ---- 摘要正文 (建议 600-1000 字, 覆盖问题/方法/结果/结论) ----
% render_paper.py 注入锚: 下一行在渲染时被替换为 abstract 节内容; 手写论文时删除锚点行直接写正文
% __ABSTRACT__

\vspace{1em}
\noindent{\heiti 关键词:} [关键词1]; [关键词2]; [关键词3]; [关键词4]

\newpage

% ============================================================
% 目录 (可选: 评审偏好不一, 不需要时保持注释状态即可)
% ============================================================
% \tableofcontents
% \newpage

% ============================================================
% 正文
% render_paper.py 注入锚: 下方成对 SECTIONS 锚行之间的
% 演示正文会被 sections/*.tex 自动替换; 手写论文时删除两个锚点行直接写
% ============================================================
% __SECTIONS__

\section{问题重述}

\subsection{问题背景}

[用两到三段交代问题背景与研究价值, 不照抄题面。]

\subsection{问题提出}

\begin{itemize}
  \item [子问题 1 的任务描述];
  \item [子问题 2 的任务描述];
  \item [子问题 3 的任务描述]。
\end{itemize}

\section{问题分析}

[逐问分析问题本质、变量关系与求解思路, 并给出总体技术路线。]

\section{模型假设与符号说明}

\subsection{模型假设}

\begin{enumerate}
  \item [假设 1 及其合理性依据];
  \item [假设 2 及其合理性依据]。
\end{enumerate}

\subsection{符号说明}

\begin{table}[H]
  \centering
  \caption{符号说明表}
  \begin{tabular}{cll}
    \toprule
    符号 & 含义 & 单位 \\
    \midrule
    $i$   & 决策变量编号 & --- \\
    $x_i$ & 第 $i$ 个决策变量 & [单位] \\
    $\lambda$ & 模型参数 & [单位] \\
    \bottomrule
  \end{tabular}
\end{table}

\section{模型建立与求解}

\subsection{[子问题 1] 的模型建立}

[模型推导过程, 示例公式如下。]

\begin{equation}
  \min \; Z = \sum_{i=1}^{n} c_i x_i
\end{equation}
\begin{equation}
  \text{s.t.}\quad \sum_{i=1}^{n} a_{ij} x_i \le b_j, \quad j = 1,2,\dots,m
\end{equation}

\subsection{模型求解与结果}

[求解算法与结果分析。]

\begin{table}[H]
  \centering
  \caption{[求解结果汇总]}
  \begin{tabular}{cccc}
    \toprule
    方案 & 指标 A & 指标 B & 目标值 \\
    \midrule
    方案 1 & [值] & [值] & [值] \\
    方案 2 & [值] & [值] & [值] \\
    方案 3 & [值] & [值] & [值] \\
    \bottomrule
  \end{tabular}
\end{table}

% 图占位示例 (实际使用时取消注释并替换路径)
\begin{figure}[H]
  \centering
  % \includegraphics[width=0.8\textwidth]{figures/result.png}
  \fbox{\parbox[c][6cm][c]{0.78\textwidth}{\centering [图占位: 技术路线图 / 结果可视化]}}
  \caption{[图题: 一句话说明该图支撑的结论]}
\end{figure}

\subsection{[子问题 2] 的模型建立与求解}

[按同样结构展开其余子问题。]

\section{灵敏度分析}

[关键参数扰动设计、结果与结论。]

\section{模型评价与推广}

\subsection{优缺点}

\begin{itemize}
  \item 优点: [优点 1]; [优点 2]。
  \item 缺点: [缺点 1]。
\end{itemize}

\subsection{模型推广}

[改进方向与应用场景。]

% __END_SECTIONS__

% ============================================================
% 参考文献: GB/T 7714 顺序编码制 (手写条目示例;
% 如使用 bib 工具, 可换成 gbt7714 宏包 + BibTeX)
% ============================================================
\begin{thebibliography}{9}
  \bibitem{journal} 作者1, 作者2. 论文题目[J]. 期刊名, 2024, 20(3): 1-10.
  \bibitem{book} 作者. 专著书名[M]. 出版地: 出版社, 2023: 50-65.
  \bibitem{thesis} 作者. 学位论文题目[D]. 城市: 学校名称, 2022.
  \bibitem{online} 作者. 网页标题[EB/OL]. (2024-01-15)[2025-06-01]. https://example.com/page.
  \bibitem{english} Author A, Author B. Paper title[J]. Journal Name, 2022, 15(2): 100-110.
\end{thebibliography}

% ============================================================
% 附录: 程序代码
% ============================================================
\appendix

\section{程序代码}

\begin{lstlisting}[language=Python, caption={求解脚本示例 solve.py}]
import numpy as np
from scipy.optimize import linprog

# model parameters
c = np.array([2.0, 3.0, 1.5])          # cost coefficients
A_ub = np.array([[1.0, 1.0, 0.0],
                 [0.0, 2.0, 1.0]])
b_ub = np.array([40.0, 50.0])

res = linprog(c, A_ub=A_ub, b_ub=b_ub,
              bounds=[(0, None)] * 3, method="highs")
print("optimal value:", res.fun)
print("solution:", res.x)
\end{lstlisting}

\end{document}
